\documentclass[Screen16to9,17pt]{foils}
\usepackage{zencurity-slides}

\selectlanguage{english}

\usepackage{tabularx}

% Reused from pentest-II-foredrag + hackerworkshop-prosa (obj 8, 11).

\begin{document}

\mytitlepage{Web Attacks II -- Tools, Labs \& Defence}
{Saint Cyber -- Cyber Akademiet, 2026}

\slide{Goals for today}
\begin{list2}
\item Use assessment tools in an authorised lab: Burp, OWASP ZAP, SQLmap, Nikto
\item Demonstrate web attacks against a controlled vulnerable application
\item Evaluate preventive controls and communicate findings
\end{list2}

\slide{Plan for today}
\begin{list2}
\item Attack proxies and scanners
\item Lab: OWASP Juice Shop
\item Defences: CSP, CSRF tokens, secure coding
\item Reporting findings
\end{list2}

\slide{Time schedule}
\begin{list2}
\item 13:15--14:45 Tools + Juice Shop -- \textbf{hands-on} (booklet)
\item 14:45--15:00 Break
\item 15:00--16:00 Defences: CSP, CSRF tokens; reporting
\item 16:00--16:15 Wrap-up
\end{list2}

\slide{Attack proxies: Zap and Burp}

\hlkimage{5cm}{webscarab_logo.png}

\begin{list1}
\item OWASP recommend Zed Attack Proxy (ZAP)\\
\link{https://www.owasp.org/index.php/OWASP_Zed_Attack_Proxy_Project}\\
\link{https://code.google.com/p/zap-extensions/}
\item I recommend buying Burp Suite \link{http://portswigger.net/burp/}\\
\link{https://pro.portswigger.net/bappstore/}
\end{list1}

\slide{Zed Attack Proxy (ZAP)}

\begin{quote}
The Zed Attack Proxy (ZAP) is an easy to use integrated penetration testing tool for finding vulnerabilities in web applications.

It is designed to be used by people with a wide range of security experience and as such is ideal for developers and functional testers who are new to penetration testing.

ZAP provides automated scanners as well as a set of tools that allow you to find security vulnerabilities manually.
\end{quote}

Source: ZAP homepage \link{https://wiki.owasp.org/index.php/OWASP_Zed_Attack_Proxy_Project} and
\url{https://www.zaproxy.org/}

\slide{Burpsuite}

\begin{quote}
Burp Suite is an integrated platform for performing security testing of web applications. Its various tools work seamlessly together to support the entire testing process, from initial mapping and analysis of an application's attack surface, through to finding and exploiting security vulnerabilities.

Burp gives you full control, letting you combine advanced manual techniques with state-of-the-art automation, to make your work faster, more effective, and more fun.
\end{quote}

Burp suite contains lots of functionality proxy, spider, scanner and can be extended with Java, Python programs

\link{http://portswigger.net/burp/}\\
\link{https://pro.portswigger.net/bappstore/}

\slide{Mini proxy: Tamper Data}

\hlkimage{16cm}{tamper-data.png}

Add-on for Firefox catch and modify data before it is sent to server\\
\link{https://addons.mozilla.org/en-US/firefox/addon/tamper-data/}

\slide{Nikto web scanner}

\hlkimage{2cm}{nikto.jpg}

\begin{quote}\small
{\bf Description}
Nikto is an Open Source (GPL) web server scanner which performs
comprehensive tests against web servers for multiple items, including
over 3200 potentially dangerous files/CGIs, versions on over 625
servers, and version specific problems on over 230 servers. Scan items
and plugins are frequently updated and can be automatically updated
(if desired).
\end{quote}

\begin{list1}
\item Quick to run, checks quite a few things. I still use and find stuff with Nikto -- and you can expand it easily
\item \verb+nikto -host 127.0.0.1 -port 8080+
\item Nikto web server scanner originally from \link{http://cirt.net/nikto2}
\end{list1}

\slide{The Exploit Database}

\hlkimage{13cm}{exploit-db.png}

Check \url{http://www.exploit-db.com/} and
\url{https://www.exploit-db.com/google-hacking-database} and use Metasploit first!

\slide{Metasploit }

%\hlkimage{}{}

\begin{quote}
The world’s most used penetration testing framework

Knowledge is power, especially when it’s shared. A collaboration between the open source community and Rapid7, Metasploit helps security teams do more than just verify vulnerabilities, manage security assessments, and improve security awareness; it empowers and arms defenders to always stay one step (or two) ahead of the game.
\end{quote}
Source: \url{https://www.metasploit.com/}

\begin{list2}
\item Has database support and integration with Nmap using the command \verb+db_nmap+
\item Includes various scanners: \url{https://www.offsec.com/metasploit-unleashed/scanner-snmp-auxiliary-modules/}
\item Metasploit Unleashed - Free Online Ethical Hacking Course\\
\url{https://www.offsec.com/metasploit-unleashed/}
\end{list2}

\slide{OWASP Juice Shop Project}

I will also use the OWASP Juice Shop Tool Project as a running example. This is an application which is modern AND designed to have security flaws.

Read more about this project at:\\
\link{https://www.owasp.org/index.php/OWASP_Juice_Shop_Project}\\
\link{https://github.com/bkimminich/juice-shop}

It is recommended to buy the Pwning OWASP Juice Shop Official companion guide to the OWASP Juice Shop from \link{https://leanpub.com/juice-shop} - suggested price USD 5.99

I run this as a docker container on Debian and it would be easy for most of you to get it running.

\slide{Brute force}

\begin{list1}
\item We call it brute force -- when testing all possibilities
\end{list1}

\begin{alltt}\small
Hydra (c) by van Hauser / THC <vh@thc.org>
Syntax: hydra [[[-l LOGIN|-L FILE] [-p PASS|-P FILE]] | [-C FILE]]
[-o FILE] [-t TASKS] [-g TASKS] [-T SERVERS] [-M FILE] [-w TIME]
[-f] [-e ns] [-s PORT] [-S] [-vV] server service [OPT]

Options:
  -S        connect via SSL
  -s PORT   if the service is on a different default port, define it here
  -l LOGIN  or -L FILE login with LOGIN name, or load several logins from FILE
  -p PASS   or -P FILE try password PASS, or load several passwords from FILE
  -e ns     additional checks, "n" for null password, "s" try login as pass
  -C FILE   colon seperated "login:pass" format, instead of -L/-P option
  -M FILE   file containing server list (parallizes attacks, see -T)
  -o FILE   write found login/password pairs to FILE instead of stdout
...
\end{alltt}

\slide{Cracking passwords -- JtR and Hashcat}


\begin{quote}
John the Ripper is a fast password cracker, currently available for
many flavors of Unix (11 are officially supported, not counting
different architectures), Windows, DOS, BeOS, and OpenVMS. Its primary
purpose is to detect weak Unix passwords.
\end{quote}

\begin{list2}
\item Hashcat is the world's fastest CPU-based password recovery tool.
\item oclHashcat-plus is a GPGPU-based multi-hash cracker using a brute-force attack (implemented as mask attack), combinator attack, dictionary attack, hybrid attack, mask attack, and rule-based attack.
\item oclHashcat-lite is a GPGPU cracker that is optimized for cracking performance. Therefore, it is limited to only doing single-hash cracking using Markov attack, Brute-Force attack and Mask attack.
\item John the Ripper password cracker old skool men stadig nyttig
\end{list2}

Source:\\
\link{https://hashcat.net/wiki/}\\
\link{http://www.openwall.com/john/}

\slide{Improving the Security of your Web Applications}

%\hlkimage{}{}

\begin{quote}

\end{quote}

Start taking a critical look on your application
\begin{list2}
\item Does it need to be \emph{on the Internet} -- otherwise add firewall rules!
\item Does it implement HTTPS/TLS?
\item What is it built using? Dependencies, frameworks, security settings
\item What is the architecture -- does it include a database, is the database locked down?
\end{list2}

\slide{Content Security Policy (CSP) for XSS Prevention}

%\hlkimage{}{}

\begin{quote}
The {\bf CSP is a security configuration tool that is supported by all major browsers.} It provides settings that a developer can take advantage of to either {\bf relax or harden security rules regarding what type of code can run inside your application.}

CSP protections come in several forms, including {\bf what external scripts can be loaded, where they can be loaded, and what DOM APIs are allowed to execute the script}.

Let’s evaluate some CSP configurations that aid in mitigating XSS risk.
\end{quote}

\begin{list2}
\item Very highly recommended, easy to check, quite easy to implement
\item Tools available for you to generate even complex CSPs
\item and lots of references:\\
\link{https://developers.google.com/web/fundamentals/security/csp}\\
\link{https://developer.mozilla.org/en-US/docs/Web/HTTP/CSP}
\end{list2}

\slide{Scanning for HTTP settings}

%\hlkimage{}{}

Nmap can also report these settings:
\begin{alltt}
|     X-XSS-Protection: 1;mode=block
|     Content-Security-Policy: script-src 'self' 'unsafe-inline' 'unsafe-eval'
|     X-Content-Type-Options: nosniff
\end{alltt}

\begin{list2}
\item If you have many sites, Nmap can also report this
\item Use the documentation from Mozilla Observatory for full explainations about these settings
\end{list2}

\slide{Better: CSRF Tokens}

\hlkimage{16cm}{was-csrf-token.png }
Source: \emph{Web Application Security}, Andrew Hoffman, 2020, ISBN: 9781492053118

\begin{list2}
\item Where do you come from, where do your browser go
\end{list2}

\slide{Example CSRF Protection Django}

%\hlkimage{}{}

\begin{quote}{\bf
Cross Site Request Forgery protection}\\
The CSRF middleware and template tag provides easy-to-use protection against Cross Site Request Forgeries. This type of attack occurs when a malicious website contains a link, a form button or some JavaScript that is intended to perform some action on your website, using the credentials of a logged-in user who visits the malicious site in their browser. A related type of attack, ‘login CSRF’, where an attacking site tricks a user’s browser into logging into a site with someone else’s credentials, is also covered.

The {\bf first defense against CSRF attacks is to ensure that GET requests (and other ‘safe’ methods, as defined by RFC 7231\#section-4.2.1) are side effect free}. Requests via ‘unsafe’ methods, such as POST, PUT, and DELETE, can then be protected by following the steps below.
\end{quote}
Source: \link{https://docs.djangoproject.com/en/4.1/ref/csrf/}

\begin{list2}
\item Many existing options can help with CSRF protection, so make sure to include this in your architecture and planning
\end{list2}

\demo{ex:nikto-webscanner}


\demo{ex:sslscan}
\demo{ex:web-site-check}

\slide{OWASP Web Security Testing Guide}

\begin{quote}
The Web Security Testing Guide (WSTG) Project produces the premier cybersecurity
testing resource for web application developers and security professionals.

The WSTG is a comprehensive guide to testing the security of web applications and
web services. Created by the collaborative efforts of cybersecurity professionals
and dedicated volunteers, the WSTG provides a framework of best practices used by
penetration testers and organizations all over the world.
\end{quote}

\begin{list2}
\item Project from OWASP:\\
\link{https://owasp.org/www-project-web-security-testing-guide/}
\item Use the Tab \emph{Release Versions} to download version 4.2 in PDF
\item Also available as a checklist \verb+OWASPv4_Checklist.xlsx+
\end{list2}

\slide{Reporting findings -- technical \& management}
% FIXME document evidence, impact, prioritised remediation (obj 11)
\begin{list2}
\item \textbf{FIXME:} document evidence, impact, prioritised remediation (obj 11)
\end{list2}

\slide{Conclusion Pentest web attacks}


\hlkimage{9cm}{eugen-str-CrhsIRY3JWY-unsplash.jpg}
~
\begin{list2}
\item You really should try testing, and investigate your existing devices
all of them
\item This is just one small part of your security posture
\item Harden servers, networks, configurations
\item Also block outgoing connections -- shellshock, log4shell, ... next CVE coming

\end{list2}

\myquestionspage


\end{document}

\slide{References}
\begin{list2}
\item \emph{Web Application Security}, Andrew Hoffman, 2020, ISBN: 9781492053118 (WAS)
\item Juice Shop: \url{https://juice-shop.github.io/} -- guide \url{https://pwning.owasp-juice.shop/companion-guide/latest/}
\end{list2}

\end{document}
